% =====================================================================
\chapter{Regla de la cadena y modelos de Markov}
\label{tema:2}
% =====================================================================

Este tema descompone la probabilidad de una secuencia en factores de la forma $P(x_t\mid x_{1:t-1})$ y presenta la aproximación de Markov que da lugar a los modelos de $n$-gramas. Los resultados de probabilidad siguen a~\cite[§1.3]{bt}; la parte de $n$-gramas, a~\cite[§3.1]{jm3}.

\section{Probabilidad condicionada}

\begin{definicion}[Probabilidad condicionada]
Sea $(\Omega,\mathcal{F},P)$ un espacio de probabilidad y sean $A,B\in\mathcal{F}$ con $P(A)>0$. La \emph{probabilidad de $B$ condicionada a $A$} es
\[
P(B\mid A)=\frac{P(A\cap B)}{P(A)}.
\]
\end{definicion}

Para $A$ fijo, la aplicación $B\mapsto P(B\mid A)$ es a su vez una medida de probabilidad sobre $(\Omega,\mathcal{F})$~\cite[§1.3]{bt}. Despejando de la definición se obtiene la \emph{regla del producto}, $P(A\cap B)=P(A)\,P(B\mid A)$.

\section{La regla de la cadena}

\begin{teorema}[Regla de la cadena]\label{teo:cadena}
Sean $A_1,\dots,A_T$ sucesos con $P(A_1\cap\dots\cap A_{T-1})>0$. Entonces
\[
P\Bigl(\,\bigcap_{t=1}^{T}A_t\Bigr)=\prod_{t=1}^{T}P\bigl(A_t\mid A_1\cap\dots\cap A_{t-1}\bigr),
\]
donde el factor con $t=1$ se interpreta como $P(A_1)$.
\end{teorema}

\begin{corolario}\label{cor:cadena-secuencias}
Para un modelo de lenguaje y una secuencia $x_{1:T}$ con $P(x_{1:T-1})>0$,
\[
P(x_{1:T})=\prod_{t=1}^{T}P(x_t\mid x_{1:t-1}).
\]
\end{corolario}

La igualdad del \cref{cor:cadena-secuencias} es exacta. Todo modelo de lenguaje autorregresivo, desde un modelo de bigramas hasta un Transformer, consiste en una forma concreta de calcular cada factor $P(x_t\mid x_{1:t-1})$.

\section{Hipótesis de Markov y modelos de \texorpdfstring{$n$}{n}-gramas}

El contexto $x_{1:t-1}$ crece con $t$, y el número de contextos posibles crece exponencialmente (\cref{prop:contar-secuencias}). La solución clásica consiste en truncar el contexto~\cite[§3.1.1]{jm3}.

\begin{definicion}[Modelo de $n$-gramas]\label{def:ngramas}
Sea $n\ge 1$. Un modelo de lenguaje satisface la \emph{hipótesis de Markov de orden $n-1$} si
\[
P(x_t\mid x_{1:t-1})=P(x_t\mid x_{t-n+1:t-1})\qquad\text{para todo } t,
\]
es decir, si la distribución del token siguiente depende solo de los $n-1$ tokens anteriores. El modelo se llama de \emph{unigramas} ($n=1$), de \emph{bigramas} ($n=2$) o de \emph{trigramas} ($n=3$).
\end{definicion}

\begin{definicion}[Matriz estocástica]
Una matriz $P\in\R^{K\times K}$ es \emph{estocástica por filas} si $P_{ij}\ge 0$ para todo $i,j$ y $\sum_{j}P_{ij}=1$ para todo $i$. Un modelo de bigramas queda determinado por la matriz estocástica $P_{ij}=P(x_t=j\mid x_{t-1}=i)$, llamada \emph{matriz de transición}.
\end{definicion}

\begin{proposicion}\label{prop:parametros}
Un modelo de $n$-gramas sobre un vocabulario de tamaño $K$ se especifica con una tabla de $K^{n}$ probabilidades, de las cuales $K^{n-1}(K-1)$ son libres.
\end{proposicion}

Con $K=5\cdot 10^{4}$ y $n=3$, la tabla tendría $1{,}25\cdot 10^{14}$ entradas. Esta explosión combinatoria es la razón por la que los modelos modernos sustituyen la tabla por una función con parámetros compartidos (módulos~9 y~13).

\addtocounter{definicion}{1}%

\addtocounter{figure}{1}%

\section{El símbolo de fin y la normalización}

La \cref{def:ml} exige que las probabilidades de \emph{todas} las secuencias finitas sumen~1. El símbolo \fin\ de fin de secuencia es imprescindible para ello~\cite[§3.1, ejercicio 3.5]{jm3}.

\begin{proposicion}\label{prop:longitud-fija}
Sea un modelo de bigramas sin símbolo de fin, con distribución inicial $\pi$ sobre $\V$ y matriz de transición estocástica $P$. Para cada longitud $L\ge 1$,
\[
\sum_{x_{1:L}\in\V^{L}}\pi(x_1)\prod_{t=2}^{L}P_{x_{t-1}x_t}=1.
\]
En consecuencia, la suma sobre todas las longitudes diverge y el modelo no define una distribución sobre $\V^{*}$.
\end{proposicion}

\begin{observacion}
Con el símbolo de fin, la generación se detiene la primera vez que aparece \fin, y cada longitud recibe solo una parte de la masa total. Si desde cualquier token se alcanza \fin\ con probabilidad positiva, la masa que queda en secuencias de longitud $L$ decrece geométricamente con $L$ y la suma total es~1; el argumento completo es el de las cadenas de Markov absorbentes~\cite[cap.~7]{bt}.
\end{observacion}

